\section*{Chapitre \ref{ch:s1:calendar-spreads-and-curve-trades} --- Calendar spreads et trades de courbe}

\begin{solution}{exo:s1:calendar-spreads-and-curve-trades:1}
$\ln 0.5/\ln 0.949 = 13.2$ jours de bourse (13.3 à partir du coefficient non arrondi).
\end{solution}

\begin{solution}{exo:s1:calendar-spreads-and-curve-trades:2}
$-0.01 \times 252/21 = -0.12$, soit un contango de 12~\% par an. Un carry de $-575.5~\%$ correspond à un spread logarithmique de $-5.755 \times 21/252 = -0.48$, de sorte que le prix proche
valait $e^{-0.48} = 0.62$ fois le prix éloigné.
\end{solution}

\begin{solution}{exo:s1:calendar-spreads-and-curve-trades:3}
Ce jour-là, le contrat~1 est un autre contrat (le contrat~2 de la veille)~: la variation de «~contrat~1 moins contrat~2~» mélange les anciennes et les nouvelles paires, et
le décalage entre les deux n'est pas un P\&L que quiconque a gagné. Le spread doit être suivi en identifiant chaque contrat.
\end{solution}

\begin{solution}{exo:s1:calendar-spreads-and-curve-trades:4}
Environ $0.30 \times 10~\% - 0.05 \times 9.9~\% = 2.5~\%$ par an, soit l'essentiel du rendement brut.
\end{solution}

\begin{solution}{exo:s1:calendar-spreads-and-curve-trades:5}
En mars, le contango s'est élargi bien au-delà de tout ce qu'avait connu la fenêtre de vingt jours de la règle~: celle-ci était longue du spread (pari sur un resserrement), tandis que le
spread continuait de s'élargir à mesure que les capacités de stockage se remplissaient. En avril, la règle est revenue avec des positions plus petites, dimensionnées sur une volatilité bien plus élevée, et une partie du mouvement du spread s'est inversée.
\end{solution}

\begin{solution}{exo:s1:calendar-spreads-and-curve-trades:6}
Ce seuil a été retenu parmi quatre, après avoir vu la seule crise de stockage qui décide du résultat~; tout seuil compris entre le contango normal et les niveaux de 2020
niveaux donnerait de bons résultats a posteriori. Une preuve supposerait un seuil fixé à l'avance à partir du mécanisme (coûts de stockage, inventaires) et testé sur
des crises qui n'ont pas servi à le choisir.
\end{solution}

\begin{solution}{exo:s1:calendar-spreads-and-curve-trades:7}
Au quotidien, les petites variations journalières du z-score font bouger la position chaque jour, et à 4 points de base par jambe et par unité négociée, le coût de rotation dépasse ce que
rapporte le signal lent ($-0.29$ après coûts contre 1.00 avant). À l'échelle mensuelle, la position change douze fois par an~; le signal, dont l'information
s'érode sur un an, perd peu à attendre, et 0.76 sur 1.05 subsiste.
\end{solution}

\begin{solution}{exo:s1:calendar-spreads-and-curve-trades:8}
Les jambes couvrent le niveau, non la pente. Quand le stockage est épuisé, la pente n'a plus de borne, et avec un effet de levier de dix, un mouvement comparable à celui de mars 2020 (58~\% du capital
à une cible de volatilité de 10~\%) représenterait un multiple du capital. Vingt ans sans crise de stockage ne sont pas une preuve qu'il n'y en a pas.
\end{solution}

\begin{solution}{pb:s1:calendar-spreads-and-curve-trades:1}
\begin{enumerate}
\item Positions sur plusieurs échéances, neutres au niveau, qui exploitent la forme de la courbe~; le retour d'un spread vers son niveau normal.
\item $d(\ln F(T_a) - \ln F(T_b)) = (T_b - T_a)\,dc$, car les termes de niveau et saisonnier s'annulent.
\item Les rendements d'opportunité baissent de façon non linéaire quand les inventaires augmentent~; la base et les rendements passés reflètent les inventaires et les primes anticipées.
\item Le coût de stockage (portage complet, ou full carry)~; il cesse de tenir lorsque le stockage est plein.
\item Par numéro de série~: la paire éloignée de son échéance est détenue et suivie d'un jour à l'autre.
\item Une autocorrélation quotidienne de 0.949, soit une demi-vie de 13.3 jours.
\item Un z-score à 20 jours, entrée au-delà de 1.5, sortie en deçà de 0.5 ou au changement de signe, dimensionné à une volatilité de 10~\%, un cent par baril par unité négociée.
\item 0.30 et 0.05.
\item Le contrat de mai a réglé à $-\$37.63$, premier prix négatif depuis 37 ans, dans un contexte de surproduction, d'effondrement de la demande et de capacités de stockage à Cushing proches
de leur capacité.
\item Le \omterm{def:s1:calendar-spreads-and-curve-trades:storage}{risque de saturation des stocks}~: que le contango s'élargisse sans borne lorsque le stockage est épuisé.
\item Position longue sur le spread depuis le 10 janvier, carry tombé à $-216~\%$ au 31 mars et $-575~\%$ le 21 avril~; $-58.3~\%$ en mars, $+22.9~\%$ en avril.
\item Pertes de 28.1~\%, 17.8~\% et 8.8~\% pour des seuils de $-80~\%$, $-40~\%$ et $-20~\%$~; le meilleur a été retenu a posteriori.
\item \textbf{Résultat nommé.} De 1986 à 2019, la règle de \omterm{def:s1:calendar-spreads-and-curve-trades:mr}{retour à la moyenne du spread} WTI affichait un ratio de Sharpe de 0.30 avant coûts et 0.05 après~; en 2020, elle a perdu
58.3~\% du capital en mars, 22.9~\% récupérés en avril, et 42.7~\% perdus sur l'année, à une cible de volatilité de 10~\%.
\item Les inventaires et les coûts de stockage, fixés avant le test.
\item Un signal lent doit être négocié lentement~: $-0.29$ par jour, 0.76 par mois après coûts.
\item La base ne prévoyait pas le VIX au comptant, mais prévoyait les rendements des contrats à terme sur le VIX~; le trade de contango couvert était rentable.
\item Un pic de volatilité qui transforme le contango en backwardation.
\item Le calendar spread de la période de roll.
\item Le chapitre~20 négociait le niveau du carry d'un marché à l'autre~; un calendar spread négocie sa variation au sein d'un même marché.
\item Il couvre le niveau de la courbe, non sa pente, et la pente n'est pas bornée lorsque le stockage est épuisé.
\end{enumerate}
\end{solution}

\begin{solution}{iq:s1:calendar-spreads-and-curve-trades:1}
Les coûts de stockage et les intérêts (la borne supérieure du contango), le rendement d'opportunité de la détention de stocks (qui augmente quand les inventaires baissent, ce qui donne
une backwardation), la saisonnalité de l'offre et de la demande, et les primes de risque versées par les hedgers.
\end{solution}

\begin{solution}{iq:s1:calendar-spreads-and-curve-trades:2}
Une série par contrat (par mois de livraison), avec ses dates d'échéance et de premier avis, plus un calendrier de roll~; les séries génériques sont des vues dérivées, jamais la
source de vérité stockée, et chaque P\&L est calculé contrat par contrat.
\end{solution}

\begin{solution}{iq:s1:calendar-spreads-and-curve-trades:3}
Vérifier les inventaires et le coût de stockage au point de livraison, les positions des \omterm{def:s1:index-rebalancing:fund}{fonds indiciels} et des ETF proches de leur roll, et les avis de la bourse~; réduire
la position si le contango s'approche ou dépasse le portage complet (full carry), plutôt que de l'alourdir sur un signal de retour à la moyenne.
\end{solution}

\begin{solution}{iq:s1:calendar-spreads-and-curve-trades:4}
Les bourses appliquent en général des marges plus faibles aux spreads qu'aux outrights, car les jambes se compensent~; le portefeuille doit être soumis à des tests de résistance sur la pente elle-même, avec des crises de
stockage historiques, un contango hypothétique au niveau du portage complet et au-delà, et la liquidité des jambes proche de l'échéance.
\end{solution}

\begin{solution}{iq:s1:calendar-spreads-and-curve-trades:5}
Les acheteurs de protection paient une prime, de sorte que les contrats à terme se situent au-dessus du VIX attendu et convergent vers le comptant~; le risque est un pic, lorsque le VIX au comptant bondit au-dessus des
contrats à terme, et le vendeur perd en quelques jours plusieurs mois de prime.
\end{solution}

\begin{solution}{iq:s1:calendar-spreads-and-curve-trades:6}
Avec $s_{t+1} = \phi s_t + \varepsilon_{t+1}$, le P\&L vaut $-s_t(s_{t+1} - s_t) = (1 - \phi)s_t^2 - s_t\varepsilon_{t+1}$, d'espérance
$(1 - \phi)\sigma^2/(1 - \phi^2) = \sigma^2/(1 + \phi)$. Sa variance est dominée par $E[s_t^2]\sigma^2 = \sigma^4/(1 - \phi^2)$ pour $\phi$ proche de un, si bien que
le ratio de Sharpe quotidien vaut environ $\sqrt{(1 - \phi)/(1 + \phi)}$. La position varie de $s_{t+1} - s_t$, soit à peu près $\varepsilon$, de sorte que les coûts prennent environ $k\,
\sigma\sqrt{2/\pi}$ par jour~: le trade exige $\sigma/(1 + \phi) > k\sqrt{2/\pi}$, soit un spread suffisamment bruité pour couvrir ses propres coûts de transaction.
\end{solution}
